\uuid{thxJ}
\titre{Division euclidienne}
\auteur{bodin ; hanani ; mzari}
\organisation{exo7}
\licence{CC BY-NC-SA 4.0}
\datecreate{2020-01-20}
\niveau{L1}
\module{Algèbre}
\chapitre{Polynôme, fraction rationnelle}
\sousChapitre{Division euclidienne}
\difficulte{1}
\theme{Arithmétique des polynômes, division euclidienne}
\exo7qcmid{lille-1.polynomes.011}

\contenu{
\begin{question}
\qclassification{105.01,105.02}
\qtags{motcle=division euclidienne}
Soient $A(X) = 2 X^4 + 3 X^3 - 8 X^2 - 2 X + 1$ et $B(X) = X^2+3X+1$. Soit $A = BQ+R$ la division euclidienne de $A$ par $B$.
\begin{answers}
    \bad{Le coefficient du monôme $X^2$ de $Q$ est $1$.}

    \bad{Le coefficient du monôme $X$ de $Q$ est $3$.}

    \bad{Le coefficient du monôme $X$ de $R$ est $2$.}

    \good{Le coefficient constant de $R$ est $2$.}
\end{answers}
\begin{explanations}
Faire le calcul !
$Q(X) = 2X^2-3X-1$, $R(X) = 4X+2$.
\end{explanations}
\end{question}
}
